% Bisection: repeatedly halve a bracket that contains a sign change, keeping
% the half that still brackets the root.  Cropped to its own bounding box so it
% can be dropped onto a Xournal++ page as an image.
%
%   latexmk -pdf bisection-pseudocode.tex
\documentclass[border=6pt,varwidth=11cm]{standalone}

\usepackage{algorithm2e}
\usepackage{amsmath}

\RestyleAlgo{plain}          % no float, no ruled box
\DontPrintSemicolon
\SetKwComment{Comment}{// }{}
\SetKwInOut{Input}{input}
\SetKwInOut{Output}{output}

\begin{document}
\begin{minipage}{\linewidth}
\begin{algorithm}[H]
  \textbf{Bisection}
  \BlankLine
  \Input{function $f$, bracket $[x_l, x_u]$ with $f(x_l) \cdot f(x_u) < 0$,
         tolerance $\varepsilon_s$}
  \Output{estimate $x_r$ of the root}
  \BlankLine
  $x_r \leftarrow x_l$\;
  \Repeat{$\varepsilon_a < \varepsilon_s$}{
    $x_{r,\text{old}} \leftarrow x_r$\;
    $x_r \leftarrow (x_l + x_u) \,/\, 2$\;
    \tcp{keep the half that still brackets the root}
    \uIf{$f(x_l) \cdot f(x_r) < 0$}{
      $x_u \leftarrow x_r$\;
    }
    \uElseIf{$f(x_l) \cdot f(x_r) > 0$}{
      $x_l \leftarrow x_r$\;
    }
    \Else{
      \Return $x_r$ \tcp*{exact root}
    }
    $\varepsilon_a \leftarrow \left| \dfrac{x_r - x_{r,\text{old}}}{x_r} \right| \cdot 100\%$\;
  }
  \Return $x_r$\;
\end{algorithm}
\end{minipage}
\end{document}
